\section{Background}\label{sec:background}

This section fixes the notation and collects the classical and quantum results used by the interface of Section~\ref{sec:interface} and by the analysis of Section~\ref{sec:advantage}.

\subsection{Maximum flow and minimum cut}\label{sec:bg-flow}

Let $G=(V,A)$ be a directed network with $n=|V|$ nodes, $m=|A|$ arcs, a source $s$, a sink $t\neq s$, and integer capacities $u_a>0$ for $a\in A$. For a node $v$, let $\delta^{+}(v)$ and $\delta^{-}(v)$ be the sets of arcs that leave and enter $v$. An \emph{$s$--$t$ flow} is a vector $f\in\mathbb{R}^{m}$ that satisfies the capacity constraints $0\le f_a\le u_a$ for all $a\in A$ and the conservation constraints
\begin{equation}\label{eq:conservation}
  \sum_{a\in\delta^{-}(v)} f_a \;=\; \sum_{a\in\delta^{+}(v)} f_a \qquad \text{for all } v\in V\setminus\{s,t\}.
\end{equation}
The value of the flow, $|f|=\sum_{a\in\delta^{+}(s)}f_a-\sum_{a\in\delta^{-}(s)}f_a$, is the net amount that leaves the source, and the maximum flow problem asks for a flow of maximum value. An \emph{$s$--$t$ cut} is a set of nodes $S\subset V$ with $s\in S$ and $t\notin S$, and its capacity is the sum of the capacities of the arcs that leave $S$. The max-flow min-cut theorem of Ford and Fulkerson~\cite{fordfulkerson1956} states that the maximum value of an $s$--$t$ flow equals the minimum capacity of an $s$--$t$ cut. The maximum flow problem can be solved in polynomial time~\cite{ahuja1993}.

Throughout the paper, an \emph{interdiction vector} $x\in\{0,1\}^{m}$ indicates the removed arcs ($x_a=1$ means that arc $a$ is interdicted), and
\begin{equation}\label{eq:F}
  \begin{aligned}
  F(x) = \max\bigl\{\, |f| :{}& f \text{ satisfies \eqref{eq:conservation} and}\\
  & 0\le f_a\le u_a(1-x_a)\ \text{for all } a\in A \,\bigr\}
  \end{aligned}
\end{equation}
is the value of the maximum $s$--$t$ flow after the interdiction. Evaluating $F(x)$ costs one maximum flow computation, which is polynomial in $n$ and $m$. We write $|x|=\sum_{a}x_a$ for the number of interdicted arcs, which is the Hamming weight of $x$.

\subsection{The maximum flow network interdiction problem}\label{sec:bg-mfnip}

\paragraph{Bilevel formulation.} Each arc has an interdiction cost $c_a\in\mathbb{Z}_{>0}$, and the interdictor has a budget $B$. The MFNIP is a Stackelberg game in which the interdictor, the leader, first removes a set of arcs within the budget, and the network user, the follower, then sends a maximum flow through the remaining network~\cite{wood1993,smith2020}. The optimal value of the problem is
\begin{equation}\label{eq:mfnip}
  \Fstar \;=\; \min_{x\in\{0,1\}^{m}} \bigl\{\, F(x) \;:\; c^{\top}x\le B \,\bigr\},
\end{equation}
where the inner problem $F(x)$ of \eqref{eq:F} is the follower's linear program.

\paragraph{Single-level reformulation.} Because the follower's problem is a linear program, it can be replaced by its dual, the minimum cut problem, and the two minimizations merge into one. Wood~\cite{wood1993} linearized the product of the cut and interdiction variables and obtained a MILP with three groups of binary variables. The variable $\alpha_i$ gives the side of node $i$ in the cut ($\alpha_i=0$ on the source side and $\alpha_i=1$ on the sink side), $\gamma_a$ indicates that arc $a$ crosses the cut and is \emph{not} interdicted, so that its capacity counts toward the cut, and $x_a$ indicates that arc $a$ is interdicted. The model is
\begin{equation}\label{eq:milp}
\begin{aligned}
  \min_{\alpha,\gamma,x}\quad & \sum_{a\in A} u_a\gamma_a \\
  \text{s.t.}\quad & \alpha_i-\alpha_j+\gamma_a+x_a \ge 0,\quad a=(i,j)\in A,\\
  & \alpha_s=0,\qquad \alpha_t=1, \\
  & c^{\top}x\le B, \\
  & \alpha\in\{0,1\}^{n},\quad \gamma,x\in\{0,1\}^{m}.
\end{aligned}
\end{equation}

\paragraph{Complexity.} Wood~\cite{wood1993} proved by reduction from the clique problem that the decision version of the MFNIP is strongly NP-complete. According to Zenklusen~\cite{zenklusen2010}, the result holds even when every interdiction cost is $1$ or $\infty$, that is, with unit costs and some arcs that cannot be interdicted.

\paragraph{Cardinality variant.} This paper studies the variant with unit costs, $c_a=1$ for every arc, and budget $B=k$. The feasible set and its size are
\begin{equation}\label{eq:Fk}
  \Fk = \bigl\{ x\in\{0,1\}^{m} : |x|\le k \bigr\},\quad N = |\Fk| = \textstyle\sum_{i=0}^{k}\binom{m}{i},
\end{equation}
and the problem is $\Fstar=\min_{x\in\Fk}F(x)$. We denote by $M=|\{x\in\Fk : F(x)=\Fstar\}|$ the number of optimal interdictions. Enumerating $\Fk$ costs exactly $N$ evaluations of $F$. We do not rely on any hardness result for the variant in which every arc can be interdicted at unit cost, because the statement of Section~\ref{sec:advantage} is a black-box result that holds in either case. Let $\kappa$ be the minimum number of arcs whose removal disconnects $s$ from $t$, which one maximum flow computation with unit capacities gives. If $k\ge\kappa$, the interdictor can cut every path and $\Fstar=0$. We call an instance \emph{non-degenerate} when $k<\kappa$. In that case every $x\in\Fk$ leaves an $s$--$t$ path with positive integer capacities, so $\Fstar\ge 1$.

\subsection{Quantum circuits, oracles, and the query model}\label{sec:bg-query}

We use the quantum circuit model~\cite{nielsen2010}. A quantum algorithm accesses a Boolean function $P:\{0,1\}^{m}\to\{0,1\}$ through an \emph{oracle}. The bit oracle acts as $U_P\ket{x}\ket{y}=\ket{x}\ket{y\oplus P(x)}$, and the phase oracle acts as $O_P\ket{x}=(-1)^{P(x)}\ket{x}$. The two are equivalent up to one auxiliary qubit, because with the target qubit in $\ket{-}=(\ket{0}-\ket{1})/\sqrt{2}$ we have $U_P\ket{x}\ket{-}=(-1)^{P(x)}\ket{x}\ket{-}$, an effect known as phase kickback~\cite{nielsen2010}.

In the \emph{query model}, the cost of an algorithm is the number of oracle applications, and all other operations are free or counted separately. The classical counterpart is an algorithm that obtains information about $F$ only by evaluating it at points of its choice, and each evaluation is one maximum flow computation. Comparing quantum queries with classical evaluations is the standard way to analyze Grover-type speedups~\cite{bbbv1997,bhmt2002}.

\subsection{Grover search and amplitude amplification}\label{sec:bg-grover}

Grover~\cite{grover1996} showed that a marked element among $N$ can be found with $\bigO(\sqrt{N})$ oracle queries, whereas a classical algorithm needs $\Theta(N)$ evaluations. Brassard, H{\o}yer, Mosca, and Tapp~\cite{bhmt2002} generalized the method to an arbitrary initial state. Let $A$ be a quantum algorithm without measurements and $\chi$ a Boolean function that marks the basis states considered good. We write $A\ket{0}=\ket{\Psi_1}+\ket{\Psi_0}$, where $\ket{\Psi_1}$ is the projection onto the good states. Let $a=\langle\Psi_1|\Psi_1\rangle$ be the probability of obtaining a good state when $A\ket{0}$ is measured, and let $\theta_a\in[0,\pi/2]$ satisfy $\sin^{2}\theta_a=a$. The amplification operator is
\begin{equation}\label{eq:Q}
  Q \;=\; -A\,S_0\,A^{-1}\,S_\chi ,
\end{equation}
where $S_\chi$ flips the sign of the good basis states (it is the phase oracle of $\chi$) and $S_0$ flips the sign of the all-zero state only. After $j$ applications of $Q$ to $A\ket{0}$, a measurement returns a good state with probability exactly
\begin{equation}\label{eq:sin2}
  \sin^{2}\bigl((2j+1)\theta_a\bigr)
\end{equation}
\cite[Section~2]{bhmt2002}. Grover's algorithm is the case $A=H^{\otimes m}$. When $A\ket{0}$ is the uniform superposition over a domain of $N$ elements of which $M$ are good, $a=M/N$. If $a$ is known, $j=\lfloor\pi/(4\theta_a)\rfloor$ iterations return a good state with probability at least $\max(1-a,a)$~\cite[Theorem~2]{bhmt2002}, which takes $\bigO(\sqrt{N/M})$ applications of $A$, of $A^{-1}$, and of the oracle. A useful special case for small demonstrations is $M/N=1/4$, for which $\theta_a=\pi/6$ and a single iteration finds a good state with certainty~\cite[Section~3.1]{bbht1998}.

\paragraph{Unknown number of solutions.} When $M$ is unknown, a fixed number of iterations can fail with high probability because \eqref{eq:sin2} oscillates in $j$. Boyer, Brassard, H{\o}yer, and Tapp (BBHT)~\cite{bbht1998} proposed an \emph{exponential search} for this case. It starts with $\mu=1$ and a growth factor $\lambda=6/5$ (any $\lambda$ strictly between $1$ and $4/3$ works). In each round it draws $j$ uniformly among the integers $0\le j<\mu$, applies $j$ Grover iterations, and measures. It stops if the result is marked and otherwise sets $\mu\leftarrow\min(\lambda\mu,\sqrt{N})$ and repeats. For $1\le M\le 3N/4$, the expected number of Grover iterations is at most $\tfrac{9}{2}\,\mu_0$ with $\mu_0=1/\sin(2\theta)<\sqrt{N/M}$, which is $\bigO(\sqrt{N/M})$~\cite[Theorem~3]{bbht1998}.

\subsection{Quantum minimum finding}\label{sec:bg-dh}

D\"urr and H{\o}yer~\cite{durr1996} use the exponential search as a subroutine to find the index of the minimum of an unsorted table $T[0..N-1]$. Theorem~1 of their paper states that the algorithm finds this index with probability at least $1/2$ in time $\bigO(\sqrt{N})$~\cite{durr1996}. In the proof, the element of rank $r$ is chosen as a threshold at some point with probability $1/r$, because the exponential search returns a marked index uniformly at random, and this bound still holds when the values are not all distinct. Every measurement is followed by a classical evaluation of the table at the measured index, either to check whether the index is marked (inside the exponential search) or to update the threshold. The original analysis does not count these evaluations, but our experiments do (Section~\ref{sec:method}).

\subsection{Dicke states}\label{sec:bg-dicke}

The Dicke state of weight $\ell$ on $m$ qubits is the uniform superposition of all basis states with exactly $\ell$ bits equal to $1$,
\begin{equation}\label{eq:dicke}
  \ket{D^{m}_{\ell}} \;=\; \binom{m}{\ell}^{-1/2}\sum_{x\in\{0,1\}^{m}:\;|x|=\ell}\ket{x}.
\end{equation}
B\"artschi and Eidenbenz~\cite{bartschi2019} constructed a unitary $U_{m,k}$ with the property
\begin{equation}\label{eq:umk}
  U_{m,k}\,\ket{0^{m-\ell}1^{\ell}} \;=\; \ket{D^{m}_{\ell}} \qquad \text{for all } \ell\le k,
\end{equation}
which they use to build $U_{m,k}$ inductively. Their circuit has $\bigO(km)$ gates and depth $\bigO(m)$ and needs no auxiliary qubit, even on a linear architecture with nearest-neighbor interactions only. By linearity, $U_{m,k}$ maps a superposition $\sum_{\ell\le k}\beta_\ell\ket{0^{m-\ell}1^{\ell}}$ to $\sum_{\ell\le k}\beta_\ell\ket{D^{m}_{\ell}}$. The authors use this fact with $k=m$ to prepare arbitrary symmetric states, preparing the input superposition with a linear ladder of controlled $R_y$ rotations~\cite[Section~4]{bartschi2019}.
